\chapter{The Black model (Black--Scholes on the forward)}
\label{cap:black}

The Black-76 model is the ``building block'' on which the whole project rests: every other model is compared with
it, implied volatilities are defined through it, and almost every exotic product uses it as an oracle in some limit.
In this chapter we derive the pricing formula, the Greeks, Black's equation, the gamma--theta link and the algorithm
that inverts the formula (implied volatility).

\section{Assumptions}
Under the $T$-forward measure the forward follows a driftless geometric Brownian motion, as in \eqref{eq:gbm}:
$\mathrm{d}F_t=\sigma F_t\dd W_t$, with $F_T=S_T$. The price of a European option is
\[
V_0=\DF\cdot\E^T\big[(\theta(F_T-K))^+\big].
\]
Inputs: $F,K>0$, $T>0$, $\DF$, $\sigma>0$. No $r$, no $q$.

Notation: $x=\ln(F/K)$ (the call is ITM if $x>0$), $s=\sigma\sqrt T$ (total volatility),
\begin{equation}
d_{1,2}=\frac{\ln(F/K)}{\sigma\sqrt T}\pm\frac{\sigma\sqrt T}{2}=\frac{x}{s}\pm\frac{s}{2},\qquad d_1-d_2=s.
\end{equation}

\section{The pricing formula, step by step}
\begin{proposizione}[Black-76 formula]
\begin{equation}\label{eq:black}
\boxed{V=\DF\,\theta\,\big[F\,\Phi(\theta d_1)-K\,\Phi(\theta d_2)\big]}
\end{equation}
\end{proposizione}

\begin{proof}[Derivation for the call ($\theta=+1$)]
\textbf{Step 1: split the payoff.} Write
\[
(F_T-K)^+=F_T\,\ind\{F_T>K\}-K\,\ind\{F_T>K\},
\]
so that
\[
\E^T[(F_T-K)^+]=\underbrace{\E^T[F_T\ind\{F_T>K\}]}_{(I)}-K\,\underbrace{\Q^T(F_T>K)}_{(II)} .
\]
\textbf{Step 2: the term (II).} From \eqref{eq:gbm-sol}, $F_T=F e^{sZ-s^2/2}$. Then
\[
F_T>K\iff sZ-\tfrac{s^2}{2}>\ln\tfrac KF\iff Z>\frac{-x+s^2/2}{s}=-d_2 .
\]
Hence $\Q^T(F_T>K)=\Q(Z>-d_2)=\Phi(d_2)$.

\textbf{Step 3: the term (I) via a change of measure.} Write $F_T=F\cdot e^{sZ-s^2/2}$ and apply \eqref{eq:tilt}
with $f(Z)=\ind\{Z>-d_2\}$:
\begin{align*}
\E\big[Fe^{sZ-s^2/2}\ind\{Z>-d_2\}\big]&=F\,\E\big[\ind\{Z+s>-d_2\}\big]\\
&=F\,\Q(Z>-d_2-s)=F\,\Phi(d_2+s)=F\,\Phi(d_1).
\end{align*}
\textbf{Step 4.} Multiply by $\DF$: $C=\DF[F\Phi(d_1)-K\Phi(d_2)]$. For the put, repeat the computation on the event
$\{F_T<K\}$, or use parity; this gives \eqref{eq:black} with $\theta=-1$.
\end{proof}

\begin{idea}[title={What $\Phi(d_1)$ and $\Phi(d_2)$ are}]
\begin{itemize}
\item $\Phi(d_2)=\Q^T(F_T>K)$ is the probability of exercise \emph{under the forward measure}, not a real-world
forecast. $\DF\,\Phi(d_2)=-\partial C/\partial K$ is the price of a \emph{digital} option (it pays 1 if $F_T>K$).
\item $\Phi(d_1)$ is the probability of the same event under the ``share measure'' (numéraire = underlying).
$\DF\,\Phi(d_1)$ is the delta with respect to the forward.
\end{itemize}
\end{idea}

\subsection{Normalized form and symmetries}
Dividing by $\DF\sqrt{FK}$ and using $\sqrt{F/K}=e^{x/2}$:
\begin{equation}\label{eq:black-norm}
V=\DF\sqrt{FK}\;b_\theta(x,s),\qquad b_\theta(x,s)=\theta\big[e^{x/2}\Phi(\theta d_1)-e^{-x/2}\Phi(\theta d_2)\big],
\end{equation}
with the symmetry $b_{-1}(x,s)=b_{+1}(-x,s)$.
The price therefore depends on only \emph{two} numbers: the moneyness $x$ and the total volatility $s$. Other
identities that are useful as tests: homogeneity $V(\lambda F,\lambda K)=\lambda V$; swap $C(F,K)=P(K,F)$; put--call symmetry
$C(F,K)=(K/F)\,P(F,F^2/K)$.

\subsection{Limits}
For $0<\sigma<\infty$: $(F-K)^+<c<F$ and $(K-F)^+<p<K$ (undiscounted prices). As $\sigma\to0$ the price tends to
$\DF$ times the intrinsic value; as $\sigma\to\infty$ the call tends to $\DF\,F$ and the put to $\DF\,K$. At the money
($x=0$) the price is exactly $C=\DF\,F\,\mathrm{erf}\big(s/(2\sqrt2)\big)$; to first order
\begin{equation}\label{eq:brenner}
C_{\text{ATM}}\approx\DF\,F\,\frac{s}{\sqrt{2\pi}}\approx0.4\,\DF\,F\,\sigma\sqrt T\qquad\text{(Brenner--Subrahmanyam)}.
\end{equation}

\section{The Greeks}
The \textbf{Greeks} are the derivatives of the price with respect to the inputs: they measure risk and tell us how
to hedge it. In the project they are derivatives with respect to the forward $F$ (with $K,T,\DF,\sigma$ fixed).

\subsection{An identity that simplifies everything}
\begin{equation}\label{eq:identita-phi}
F\,\varphi(d_1)=K\,\varphi(d_2).
\end{equation}
\emph{Proof.} Since $d_1-d_2=s$ and $d_1+d_2=2x/s$,
\[
\frac{\varphi(d_2)}{\varphi(d_1)}=e^{(d_1^2-d_2^2)/2}=e^{(d_1-d_2)(d_1+d_2)/2}=e^{s\cdot(2x/s)/2}=e^{x}=\frac FK .
\]

\emph{Consequence.} When \eqref{eq:black} is differentiated with respect to an input $y$, the terms arising from
$\partial_yd_1$ and $\partial_yd_2$ combine into $\DF\,F\varphi(d_1)\,\partial_y(d_1-d_2)=\DF\,F\varphi(d_1)\,\partial_ys$.
So \emph{only the change in $s$ matters}. If $y=F$, $s$ does not change and those terms vanish.

\subsection{The formulas}
\begin{align}
\Delta_F&=\frac{\partial V}{\partial F}=\theta\,\DF\,\Phi(\theta d_1),\label{eq:delta}\\
\Gamma_F&=\frac{\partial^2V}{\partial F^2}=\frac{\DF\,\varphi(d_1)}{F\,\sigma\sqrt T}\quad\text{(same for call and put)},\label{eq:gamma}\\
\vega&=\frac{\partial V}{\partial\sigma}=\DF\,F\,\varphi(d_1)\sqrt T\quad\text{(same for call and put)},\label{eq:vega}\\
\Theta&=\Theta_0+rV,\qquad\Theta_0=-\frac{\DF\,F\,\varphi(d_1)\,\sigma}{2\sqrt T},\qquad r=-\frac{\ln\DF}{T}.\label{eq:theta}
\end{align}
Derivations: for the delta, $y=F$ and the $\partial d$ terms vanish; for the gamma, $\partial_F\Phi(d_1)=\varphi(d_1)/(Fs)$;
for the vega, $\partial_\sigma s=\sqrt T$; for $\Theta_0$, $\partial_Ts=\sigma/(2\sqrt T)$ at fixed $\DF$. The vega is
``per unit of volatility'': per ``vol point'' divide by 100 (on Hull's Greeks example in the next section the vega is
12.1055, that is, 0.121 per vol point).

Check identities: $\vega=\sigma T F^2\Gamma_F$ (long gamma $\Leftrightarrow$ long vega for a single option);
$\Delta_C-\Delta_P=\DF$; $\partial C/\partial K=-\DF\,\Phi(d_2)$; $\partial^2C/\partial K^2=\DF\,\varphi(d_2)/(Ks)$ (the
lognormal density, consistent with Breeden--Litzenberger).

\subsection{Theta and its convention}
The project's theta is $\Theta=\partial V/\partial t$ in calendar time, per year (divide by 365 for a daily theta).
It holds fixed $F$ (the forward for the same expiry date), $K$, $\sigma$ and the zero rate $r=-\ln\DF/T$. So, as time
passes, the discount factor $\DF=e^{-r(T-t)}$ ``grows'' towards 1 and contributes $+rV$: this is the return on the
premium paid. The part at fixed $\DF$, $\Theta_0$, is pure time decay.

\begin{perche}[title={Why this convention?}]
It is exactly the $\mathrm{d}t$ term of Black's equation (next section): under $\Q^T$ the forward has no drift, so
this is the theta that enters the expected P\&L of a hedged option. The $rV$ term is offset by the cost of funding
the premium, and $\Theta_0$ remains.
\end{perche}

\begin{attenzione}[title={The project's theta is not Hull's}]
For a deep ITM call the theta can be \emph{positive}: $rV$ exceeds $|\Theta_0|$. In the project this happens at
50.6\% of the ITM points of the test grid. Only $\Theta_0$ is always negative; on the OTM options of the grid the
ratio $rV/|\Theta_0|$ reaches at most 0.688, so there $\Theta<0$ always.

Hull, by contrast, holds spot, $r$ and $q$ fixed. The link is derived step by step. With a continuous-yield model
$F=S\,e^{(r-q)(T-t)}$; at fixed $S$ the forward moves in time with $\partial F/\partial t=-(r-q)F$. By the chain
rule
\begin{equation}\label{eq:theta-spot}
\Theta_S=\Theta+\Delta_F\,\frac{\partial F}{\partial t}=\Theta-(r-q)\,F\,\Delta_F .
\end{equation}
Likewise $\Delta_S=\Delta_F\,\partial F/\partial S=e^{-qT}\Phi(d_1)$ for the call. On Hull's Greeks example
($S=49$, $K=50$, $r=5\%$, $q=0$, $\sigma=20\%$, $T=20/52$) the kernel theta is $-3.027$/year;
\eqref{eq:theta-spot} gives $-4.30533$/year, that is, $-0.0117954$ per day, which is Hull's value ($-4.31$/year).
\end{attenzione}

\section{Black's equation and the gamma--theta link}
By Feynman--Kac \eqref{eq:feynman-kac}, the undiscounted price $v=V/\DF$ satisfies
$\partial_tv+\frac12\sigma^2F^2v_{FF}=0$. Multiplying by $\DF$ (and recalling that $\partial_t\DF=r\DF$):
\begin{equation}\label{eq:pde-black}
\Theta+\tfrac12\sigma^2F^2\Gamma_F-rV=0 .
\end{equation}
\subsection{The P\&L of a delta-hedged option}
I buy an option and sell $\Delta_F$ forwards (or futures) to neutralize the risk to first order. Over an interval
$\Delta t$, by a Taylor expansion:
\[
\mathrm{d}\Pi\approx\underbrace{\Delta_F\,\Delta F-\Delta_F\,\Delta F}_{=0\ \text{(hedge)}}+\tfrac12\Gamma_F(\Delta F)^2+\Theta_0\,\Delta t .
\]
Using \eqref{eq:pde-black} ($\Theta_0=-\frac12\sigma^2F^2\Gamma_F$, once the cost of funding the premium $rV$ is
removed):
\begin{equation}\label{eq:pnl-gamma}
\boxed{\mathrm{d}\Pi=\tfrac12\Gamma_F F^2\Big[\Big(\frac{\Delta F}{F}\Big)^2-\sigma^2\Delta t\Big]}
\end{equation}
\begin{idea}[title={Key idea: ``gamma pays for theta''}]
Whoever is long a hedged option makes money if the \emph{realized} variance $(\Delta F/F)^2$ exceeds the
\emph{implied} one $\sigma^2\Delta t$, and loses otherwise. Theta is the ``rent'' paid for owning gamma. The gain is
weighted by gamma, so it is path-dependent: a large move far from the strike (where gamma is small) earns little.
For short-dated ATM options gamma and theta grow like $1/\sqrt T$ while vega decreases like $\sqrt T$: a portfolio
long the 1-year option and short the 1-month option is short gamma and long vega.
\end{idea}

\section{Implied volatility}
\subsection{The problem}
Given a market price $V^m$, find $\sigma$ such that $V_{\text{Black}}(\sigma)=V^m$. The problem is solved in
normalized form. The steps:
\begin{enumerate}
\item Undiscounted price $v=V/\DF$; if $v$ does not lie strictly between the intrinsic value $\iota$ and $F$ (call)
or $K$ (put), no volatility exists: \code{NaN} is returned with a status code.
\item Always switch to the \textbf{OTM} option (thanks to parity): $\beta=(v-\iota)/\sqrt{FK}$.
\item Solve $b(s):=b_{+1}(-|x|,s)=\beta$ for the unknown $s$, and finally $\sigma=s/\sqrt T$.
\end{enumerate}

\subsection{Existence and uniqueness}
Differentiating the normalized form and using \eqref{eq:identita-phi}:
\begin{equation}\label{eq:vega-norm}
b'(s)=\frac{1}{\sqrt{2\pi}}\exp\!\Big(-\frac{x^2}{2s^2}-\frac{s^2}{8}\Big)>0,\qquad
b''(s)=b'(s)\Big(\frac{x^2}{s^3}-\frac s4\Big).
\end{equation}
$b$ is strictly increasing, starts from $0$ (as $s\to0$, the OTM option is worth zero) and tends to $e^{-|x|/2}$:
every admissible $\beta$ has \emph{exactly one} solution. The true vega is $\vega=\DF\sqrt{FK}\sqrt T\,b'(s)$.

\subsection{Newton without surprises (Manaster--Koehler)}
Newton's method $s_{n+1}=s_n-(b(s_n)-\beta)/b'(s_n)$ converges quickly but can diverge from a bad start. From
\eqref{eq:vega-norm}: $b''=0$ at $s_c=\sqrt{2|x|}$; $b$ is \textbf{convex} on $(0,s_c)$ and \textbf{concave} on
$(s_c,\infty)$, and vega is maximal at $s_c$.
\begin{perche}[title={Why start from the inflection point}]
\begin{itemize}
\item If the root lies in the concave part ($\beta>b(s_c)$), the tangents lie \emph{above} the graph: every Newton
step lands \emph{before} the root, and the iterates increase monotonically towards it.
\item If the root lies in the convex part ($\beta<b(s_c)$), the tangents lie \emph{below}: the iterates decrease
monotonically towards the root.
\end{itemize}
In both cases the root is never ``overshot''. At the money ($x=0$) $s_c=0$, and the exact inverse
$s=2\sqrt2\,\mathrm{erfinv}(\beta)$ is used instead.
\end{perche}
The project's solver adds a \emph{safeguard} anyway: it keeps a bracket $[\ell,h]$ containing the root (initially
$[0,10\sqrt T]$) and, if Newton leaves the bracket, it takes a bisection step. Elements that have not converged after
100 iterations go to Brent's method (which was never needed on the test grid).

\subsection{Why always invert the OTM option}
By parity, the undiscounted OTM price equals the time value of the ITM option at the same strike, and the vega is the
same. But an ITM price carries an absolute rounding error $\approx\varepsilon\max(F,K)$, which ends up \emph{entirely}
on the time value. The error on the volatility is
\begin{equation}\label{eq:errore-iv}
|\delta\sigma|\approx\frac{|\delta v_{\text{time value}}|}{\sqrt{FK}\sqrt T\,b'(s)} .
\end{equation}
For an OTM input $\delta v\approx\varepsilon v$, for an ITM input $\delta v\approx\varepsilon\max(F,K)$: the a priori
estimate is $\lesssim10^{-14}$ against $\sim10^{-12}$.

Measured in the project, on the test grid ($147{,}600$ points per option type): round-trip error at most
${1.07\times10^{-14}}$ for OTM inputs and ${8.23\times10^{-13}}$ for ITM inputs, against a threshold of $10^{-10}$.
An independent check on $300{,}000$ random off-grid points gave ${1.05\times10^{-14}}$ and ${6.54\times10^{-13}}$. The
measurements confirm the error model.

In the extreme wings the vega vanishes: tiny errors on the price become huge errors on the volatility, and below
machine precision there is no information left.

\section{Checking the Greeks with finite differences}
The Greek formulas are checked by comparing them with numerical derivatives:
\begin{equation}\label{eq:diff-finite}
D_1f=\frac{f(u+h)-f(u-h)}{2h},\qquad D_2f=\frac{f(u+h)-2f(u)+f(u-h)}{h^2}.
\end{equation}

\begin{perche}[title={How to choose the step $h$}]
Two errors compete.
\begin{itemize}
\item The truncation (Taylor) error is $\approx h^2|f'''|/6$ for $D_1$ and $\approx h^2|f''''|/12$ for $D_2$: it grows
with $h$.
\item The rounding error is $\approx\varepsilon|f|/h$ for $D_1$ and $\approx4\varepsilon|f|/h^2$ for $D_2$: it grows
as $h$ decreases.
\end{itemize}
Minimizing the sum gives $h\propto\varepsilon^{1/3}$ for $D_1$ and $h\propto\varepsilon^{1/4}$ for $D_2$. An $h$ that
is too small is worse than one that is too large. The steps are expressed in the natural scale of the variable: for
$F$ a multiple of $F\sigma\sqrt T$ is used, which is one standard deviation of the forward.
\end{perche}

\paragraph{Theta by finite differences.} For the theta, $T$ and $\DF$ must be moved together at fixed $r$:
\begin{equation}\label{eq:theta-fd}
\Theta_{\text{FD}}=\frac{V(F,K,T-h,\DF\,e^{+rh},\sigma)-V(F,K,T+h,\DF\,e^{-rh},\sigma)}{2h}.
\end{equation}
Moving only $T$ measures $\Theta_0$ and is off by $rV$.

\begin{progetto}[title={In the qpricer project: the gamma that ``would not pass''}]
\textbf{The symptom.} With the step originally planned, $h=10^{-4}F\sigma\sqrt T$, the finite-difference gamma was off
by up to ${1.773\times10^{-5}}$ in relative terms (threshold $10^{-5}$), at 78 points out of $147{,}600$, all with
$\sigma\sqrt T\approx0.007$.

\textbf{The diagnosis.} The analytic gamma was right: the identity $\vega=\sigma TF^2\Gamma$ held below $10^{-12}$,
and larger steps converged to it. It was \emph{the reference that was noisy}. At the worst point the price noise was
about $600\,\varepsilon V$, not $\varepsilon V$, because of the cancellation in $F\Phi(d_1)-K\Phi(d_2)$. The rounding
term $4\varepsilon|f|/h^2$ was therefore 600 times larger than expected, and the optimal step grows by
$600^{1/4}\approx5$.

\textbf{The new step, step by step.} Write $h=c\,F\sigma\sqrt T$ and model the worst error on the grid as
$A c^2+B/c^2$ (truncation plus rounding).
\begin{enumerate}
\item $A=4.47$ is measured at the corner where truncation dominates ($\sigma=1$, $T=5$, $z=3$, with
$c={2\times10^{-3}}$).
\item $B={1.77\times10^{-13}}$ is measured at the corner where rounding dominates (with $c=10^{-4}$).
\item The minimum of $Ac^2+B/c^2$ is at $c^*=(B/A)^{1/4}={4.46\times10^{-4}}$, where it equals
$2\sqrt{AB}\approx{1.8\times10^{-6}}$.
\end{enumerate}
With $c={4.5\times10^{-4}}$ the measured error is ${9.96\times10^{-7}}$. The threshold was not loosened and the grid
was not changed: the reference changed, not the criterion. An independent check confirmed the analytic gamma to
${8.9\times10^{-15}}$ against finite differences computed with 60 decimal digits, and also checked the gamma of the
ITM options (${1.24\times10^{-6}}$ with its own Richardson scheme).
\end{progetto}

\section{Main numerical pitfalls}
\begin{enumerate}
\item \textbf{Normal tails.} Never $1-\Phi(d)$; use the form with $\theta$, \code{ndtr(theta*d)}.
\item \textbf{ITM cancellation.} The time value of an ITM option, the finite-difference gamma of an ITM option and
the implied volatility from an ITM price all have an absolute error $\approx\varepsilon(F+K)$. One works on the OTM
side. A measured example: for an ITM quote with $3<|z|<7.5$, where $z=\ln(K/F)/(\sigma\sqrt T)$, the solver returns a
volatility with status ``converged'' but with an error that can reach ${2.0\times10^{-4}}$ in absolute terms
($\sigma=1$, $T=5$, $z=6.5$). This is not a defect of the solver: the error matches the conditioning limit
$\varepsilon\max(v,F,K)/(\vega/\DF)$ of \eqref{eq:errore-iv}. It is, however, an unflagged error, a known and
unresolved limitation; the calibration is not affected, because it uses only OTM quotes (Chapter~\ref{cap:dati}).
\item \textbf{Parity to $10^{-12}$.} Computing call and put separately, the error is
$\approx c\,\varepsilon\,\DF(F+K)$ with $c\approx2$--4. For strikes very large relative to $F=100$ (up to
$K\approx{8.2\times10^4}$ in the test domain) the ulp of $K$ already exceeds $10^{-12}$, and that threshold is not
representable in double precision. This is why parity is checked to $10^{-12}$ for $K/F\in[0.25,4]$, and elsewhere
at the limit of double precision, $4$ ulp of $\DF\max(F,K)$. Computing the put as $C-\DF(F-K)$ would turn the test
into a tautology.
\item \textbf{$s\to0$.} $d_{1,2}=0/0$ when $F=K$. The code substitutes the limits ($+\infty$ for $x>0$, $-\infty$
for $x<0$, $0$ for $x=0$), and the formulas automatically return the correct limits, with no special branches. At the
money the vega at $\sigma=0$ is the right derivative $\DF\,F\sqrt T/\sqrt{2\pi}$, consistent with
\eqref{eq:brenner}.
\item \textbf{Deep wings.} \code{ndtr} returns 0 below $d\approx-37.5$, even though $\Phi$ would still be
representable down to $d\approx-38.5$. If $\Phi(d_2)$ underflows and $\Phi(d_1)$ does not, the normalized function
$b(s)$ loses its second term and comes out up to about 17 times too large: the solver converges on a wrong function.
An independent check showed that, for OTM prices below about $10^{-280}$, this produced a wrong volatility, even
$\sigma=0$, without flagging it. The fix: the root is accepted only if $s>0$ and
$\Phi(d_2)\ge10^3\cdot\text{(smallest normal number)}\approx{2.2\times10^{-305}}$; otherwise the result is \code{NaN}
with a dedicated status and a warning. No quoted price lies in that region, but the rule ``never a wrong number
silently'' applies everywhere.
\item \textbf{Impossible prices.} If the price does not lie strictly between the intrinsic value and $F$ (call) or
$K$ (put), or the volatility exceeds the cap $\sigma_{\max}=10$, the solver returns \code{NaN} with a status code and a
warning. The status is per element: a single stale or crossed quote does not stop the inversion of an entire chain,
and it is never silently ``adjusted''. Invalid market inputs ($F\le0$, $T\le0$, \dots), by contrast, raise an error:
they are programming errors, not data.
\end{enumerate}

\begin{progetto}[title={In the qpricer project: Phase 1 results}]
The numbers come from \code{scripts/phase1\_gates.py} and from an independent check, which uses a Black formula
computed with 60 decimal digits, the complex-step method and its own bisection solver.
\begin{center}
\small
\begin{tabularx}{\textwidth}{X l l}
\toprule
check & measured & threshold\\
\midrule
parity, $K/F\in[0.25,4]$, ${2\times10^5}$ points & ${1.137\times10^{-13}}$ & $10^{-12}$\\
parity, whole test domain & 2.25 ulp (${1.455\times10^{-11}}$) & 4 ulp\\
implied volatility, round trip, OTM input & ${1.07\times10^{-14}}$ & $10^{-10}$\\
implied volatility, round trip, ITM input & ${8.23\times10^{-13}}$ & $10^{-10}$\\
delta vs finite differences (relative) & ${2.05\times10^{-8}}$ & $10^{-5}$\\
gamma vs finite differences (relative) & ${9.96\times10^{-7}}$ & $10^{-5}$\\
vega vs finite differences (relative) & ${2.96\times10^{-7}}$ & $10^{-5}$\\
theta vs finite differences (relative) & ${1.69\times10^{-8}}$ & $10^{-5}$\\
Greeks vs 60-digit reference & $\le{5.1\times10^{-14}}$ & ---\\
Hull's example (Black-76, $F=K=20$) & $1.1166414\ldots$ vs $1.1166$ & $\pm{0.5\times10^{-4}}$\\
\bottomrule
\end{tabularx}
\end{center}
On $2\times147{,}600$ elements the Newton solver uses at most 12 evaluations (6.95 on average), and the Brent fallback
is never called. The file \code{tests/golden/bsm\_black76.json} freezes 12 cases (price and Greeks) computed in
extended precision; the implementation reproduces them to within $10^{-11}$ relative.
\end{progetto}

\begin{riassunto}
\begin{itemize}
\item $V=\DF\,\theta\,[F\,\Phi(\theta d_1)-K\,\Phi(\theta d_2)]$; it depends only on $x=\ln(F/K)$ and $s=\sigma\sqrt T$.
\item $F\,\varphi(d_1)=K\,\varphi(d_2)$ simplifies all the Greeks. $\vega=\sigma TF^2\Gamma$.
\item Theta at fixed $F$ and $r$: $\Theta=\Theta_0+rV$, which can be positive for ITM options; Hull's differs by
$(r-q)F\Delta_F$.
\item Black's PDE $\Rightarrow$ hedged P\&L $=\frac12\Gamma F^2[(\Delta F/F)^2-\sigma^2\Delta t]$.
\item Implied volatility is unique; Newton started from the point of maximum vega converges monotonically; invert the
OTM option.
\item Finite differences: the step balances truncation and rounding, and must be chosen on the true price noise.
\end{itemize}
\end{riassunto}

\begin{domande}
\begin{enumerate}
\item Derive Black's formula for the call using the change of measure. Where does $\Phi(d_1)$ come in?
\item You are long an ATM call, delta-hedged every day. What is your P\&L? \emph{Hint:} \eqref{eq:pnl-gamma}.
\item Why does Newton start from $s_c=\sqrt{2|x|}$ for the implied volatility?
\item The theta of an ITM call is positive: is that a mistake? \emph{Hint:} no, $\Theta=\Theta_0+rV$ and $rV$ can
exceed $|\Theta_0|$.
\item How does one go from the kernel theta to Hull's? \emph{Hint:} \eqref{eq:theta-spot}.
\item The finite-difference gamma fails the test: should the threshold be loosened? \emph{Hint:} no; analyze the
noise of the reference and choose the optimal step.
\end{enumerate}
\end{domande}
